\chapter{Conclusion and Future Work}
\label{ch:conclusion}

\section{Conclusion}
This project built and evaluated hybrid quantum--classical neural networks
for detecting tuberculosis in chest radiographs, in two placements: a \cq{}
model in which a CNN's features are classified by a four-qubit variational
circuit, and a \qc{} model in which a four-qubit quanvolution filter
transforms image patches before a CNN. Each hybrid was compared with
classical controls of the same structure and size, with three seeds, five
training-set sizes and paired McNemar tests, and the circuits were analysed
directly. The main findings are:
\begin{enumerate}
  \item The \cq{} hybrid reached $98.41 \pm 0.32\%$ accuracy (F1
        $0.952 \pm 0.010$) against $97.72 \pm 0.40\%$ for a matched
        classical head and $97.67 \pm 0.33\%$ for the CNN; the gain was
        consistent across seeds but not significant ($p = 0.09$--$1.00$).
        With ResNet18 the quantum head reached $99.47 \pm 0.18\%$, and the
        backbone mattered about three times more than the head.
  \item The \qc{} hybrid matched the CNN with four times fewer parameters,
        but so did average pooling and a learned $2\times2$ convolution: the
        efficiency comes from downsampling, not from the quantum filter.
  \item With 2\% of the data the \qc{} hybrid was the most stable model;
        bottleneck heads were unstable, although the quantum head still
        ranked images best (ROC-AUC 0.927).
  \item Entanglement made the circuit 20--60 times more expressive but
        changed accuracy by at most one test image; angle encoding was
        consistently better than amplitude encoding.
  \item Gradient variance halved with each added qubit, supporting small
        circuits, and the trained model kept its accuracy with 10 shots and
        depolarising noise up to $p = 0.2$, with exact parameter-shift
        gradients, so it is directly transferable to NISQ hardware.
\end{enumerate}
Small quantum circuits can therefore be inserted into a medical image
classifier without loss and, in the \cq{} placement, with a small,
consistent gain, but at the scale that can be simulated today they give no
statistically demonstrable advantage over an equally sized classical
component.

\section{Limitations}
All circuits were simulated with 2 to 10 qubits, and noise was modelled by
an idealised depolarising channel. Images were reduced to $64 \times 64$
pixels, which may hide subtle findings. All results come from one public
dataset whose normal and TB images were partly collected from different
sources, so models may exploit acquisition differences such as markers and
borders (Figure~\ref{fig:quanvmaps}), and performance on images from another
hospital could be lower. With 630 test images and three seeds, differences
below about one percentage point cannot be confirmed. The quanvolution
filter was random and frozen, hyperparameters were shared rather than tuned
per model, and training times were measured on a shared laptop.

\section{Future Work}
The trained \cq{} head should be run on a cloud quantum processor with the
parameter-shift rule, readout error mitigation and zero-noise
extrapolation~\cite{temme2017}. The quanvolution filter could be made
trainable and deeper so that its entanglement can contribute, and data
re-uploading circuits~\cite{perezsalinas2020} and quantum kernel
methods~\cite{havlicek2019} should be compared with the present heads. The
low-data collapse of bottleneck heads could be addressed by tuning the
decision threshold on validation data or by a class-weighted loss. Training
at $224 \times 224$ pixels and testing on independent collections such as
the Montgomery and Shenzhen datasets~\cite{jaeger2014} would measure
generalisation across sources, and saliency methods such as
Grad-CAM~\cite{selvaraju2017} would show whether the models attend to the
lungs rather than to markers, as Rahman et al.~\cite{rahman2020} did for
classical models. More seeds and cross-validation would provide the
statistical power needed to confirm or reject the small advantage of the
\cq{} head.
